%%═══════════════════════════════════════════════════════════════════
%% Lecture 15 — PDEs: Classification and Elliptic Equations
%% 77315 — Computational Physics
%%═══════════════════════════════════════════════════════════════════

\section{Lecture 15 — PDE Classification \& Elliptic Equations}\label{sec:lec15}
\hrule\vspace{1.5em}

%%──────────────────────────────────────────────────────────────────
\subsection*{Overview}
We classify second-order PDEs and focus on \emph{elliptic} equations
(Poisson/Laplace).  Finite-difference discretisation converts the PDE
into a large sparse linear system.  Three iterative solvers are
discussed: Jacobi, Gauss–Seidel, and Successive Over-Relaxation (SOR).

%%──────────────────────────────────────────────────────────────────
\subsection{Classification of Second-Order PDEs}\label{subsec:lec15:class}

A general 2nd-order PDE in two variables $u(x,y)$:
\begin{equation}
  A\,u_{xx} + B\,u_{xy} + C\,u_{yy} + \ldots = 0.
  \label{eq:lec15:pde}
\end{equation}
The discriminant $\Delta = B^2 - 4AC$ determines the type:

\begin{center}
\begin{tabular}{lll}
\toprule
$\Delta$ & Type & Example \\
\midrule
$>0$ & Hyperbolic & Wave equation $u_{tt} = c^2 u_{xx}$ \\
$=0$ & Parabolic  & Diffusion $u_t = D\,u_{xx}$ \\
$<0$ & Elliptic   & Poisson $u_{xx}+u_{yy}=f$ \\
\bottomrule
\end{tabular}
\end{center}

Elliptic problems are global BVPs (no time direction); hyperbolic and
parabolic problems evolve in time.

%%──────────────────────────────────────────────────────────────────
\subsection{Poisson Equation and Finite Differences}\label{subsec:lec15:poisson}

The 2D Poisson equation on a grid with $N$ interior points per side:
\begin{equation}
  \nabla^2 u = \frac{\partial^2 u}{\partial x^2}
              +\frac{\partial^2 u}{\partial y^2} = f(x,y).
  \label{eq:lec15:poisson}
\end{equation}
Discretise with spacing $h = 1/(N+1)$ using 5-point stencil:
\begin{equation}
  \frac{u_{i+1,j}+u_{i-1,j}+u_{i,j+1}+u_{i,j-1}-4u_{i,j}}{h^2} = f_{ij}.
  \label{eq:lec15:stencil}
\end{equation}
This generates an $N^2\times N^2$ sparse linear system
$A\mathbf{u}=\mathbf{b}$ with at most 5 non-zero entries per row.
Direct solvers require $\bigO{N^4}$ time; iterative methods are preferred.

%%──────────────────────────────────────────────────────────────────
\subsection{Jacobi Iteration}\label{subsec:lec15:jacobi}

Rearrange~\eqref{eq:lec15:stencil} to update $u_{ij}$ from neighbours:
\begin{equation}
  u_{ij}^{\rm new} = \tfrac{1}{4}\!\left(
    u_{i+1,j}+u_{i-1,j}+u_{i,j+1}+u_{i,j-1}-h^2 f_{ij}
  \right).
  \label{eq:lec15:jacobi}
\end{equation}
All updates use old values.  Convergence to precision $10^{-p}$ requires
approximately
\begin{equation}
  n_{\rm iter} \approx \tfrac{1}{2}\,p\,N^2
  \label{eq:lec15:jacobi_iters}
\end{equation}
iterations — very slow for large $N$.

%%──────────────────────────────────────────────────────────────────
\subsection{Gauss--Seidel Iteration}\label{subsec:lec15:gs}

Same as Jacobi but use updated values \emph{immediately}:
\begin{equation}
  u_{ij}^{\rm new} = \tfrac{1}{4}\!\left(
    u_{i+1,j}^{\rm old}+u_{i-1,j}^{\rm new}+u_{i,j+1}^{\rm old}
    +u_{i,j-1}^{\rm new}-h^2 f_{ij}
  \right).
  \label{eq:lec15:gs}
\end{equation}
Gauss–Seidel converges about \textbf{twice as fast} as Jacobi and
requires no extra storage.

%%──────────────────────────────────────────────────────────────────
\subsection{Successive Over-Relaxation (SOR)}\label{subsec:lec15:sor}

Introduce a relaxation parameter $\omega\in(1,2)$ to accelerate
Gauss–Seidel:
\begin{equation}
  u_{ij}^{\rm new} = (1-\omega)\,u_{ij}^{\rm old}
                   + \omega\,u_{ij}^{\rm GS},
  \label{eq:lec15:sor}
\end{equation}
where $u_{ij}^{\rm GS}$ is the Gauss–Seidel update.  The optimal
relaxation parameter for the Poisson equation on an $N\times N$ grid is
\begin{equation}
  \omega_{\rm opt} = \frac{2}{1+\sin(\pi/(N+1))}.
  \label{eq:lec15:omega_opt}
\end{equation}
With optimal $\omega$, SOR converges in $\bigO{N}$ iterations rather
than $\bigO{N^2}$ — a spectacular speedup for large grids.

\nt{$\omega=1$ recovers Gauss–Seidel.  $\omega\ge 2$ is unconditionally divergent.  Choosing $\omega$ slightly above 1 is always safe if $\omega_{\rm opt}$ is unknown.}
